\chapter{Carga Final a SQLite}

La persistencia final del proyecto se materializa en \texttt{db/lab1\_desigualdad.sqlite}. En términos operativos, esta elección deja un artefacto local único, consultable con SQL y suficiente para validar consistencia entre dimensión, tabla de hechos y metadata sin depender de un servidor externo.

\section{Estructura y conteos}

La base resultante contiene las tablas \texttt{dim\_comuna}, \texttt{fact\_desigualdad\_comunal} y \texttt{metadata\_fuentes}. El reporte de carga y la inspección directa de la base coinciden en los conteos observados.

\begin{table}[htbp]
\centering
\footnotesize
\caption{Conteo de filas por tabla en la base SQLite final.}
\begin{tabular}{|p{4.2cm}|r|p{6.0cm}|}
\hline
Tabla & Filas & Rol \\
\hline
dim\_comuna & 32 & Dimensión comunal \\
fact\_desigualdad\_comunal & 32 & Tabla de hechos \\
metadata\_fuentes & 4 & Contrato operativo de fuentes \\
\hline
\end{tabular}
\end{table}


\section{Integridad básica}

El archivo abre correctamente y el chequeo \texttt{PRAGMA quick\_check} devuelve \texttt{ok}. Además, los reportes de carga informan ausencia de duplicados por \texttt{codigo\_comuna} en la dimensión y en la tabla de hechos, y ausencia de códigos huérfanos entre ambas tablas.

\section{Consultas de prueba}

Las consultas documentadas tras la carga reproducen conteos por tabla y rankings básicos para áreas verdes, pobreza e IPP. Esto confirma que la base queda utilizable no solo como respaldo del dataset final, sino también como soporte de consulta analítica.

\begin{table}[htbp]
\centering
\scriptsize
\caption{Consultas de prueba documentadas tras la carga en SQLite.}
\begin{tabular}{|p{3.2cm}|p{0.9cm}|p{3.0cm}|p{1.2cm}|p{2.8cm}|p{2.3cm}|p{1.8cm}|}
\hline
Consulta & Ord. & Tabla & Código & Comuna & Valor & Detalle \\
\hline
conteo\_filas\_por\_tabla & - & dim\_comuna & - & - & 32.0 & conteo de filas \\
conteo\_filas\_por\_tabla & - & fact\_desigualdad\_comunal & - & - & 32.0 & conteo de filas \\
conteo\_filas\_por\_tabla & - & metadata\_fuentes & - & - & 4.0 & conteo de filas \\
top\_5\_menor\_areas\_verdes\_m2\_hab & 1.0 & fact\_desigualdad\_comunal & 13104 & CONCHALI & 0.603247 & areas\_verdes\_m2\_hab \\
top\_5\_menor\_areas\_verdes\_m2\_hab & 2.0 & fact\_desigualdad\_comunal & 13102 & CERRILLOS & 0.878035 & areas\_verdes\_m2\_hab \\
top\_5\_menor\_areas\_verdes\_m2\_hab & 3.0 & fact\_desigualdad\_comunal & 13108 & INDEPENDENCIA & 0.878941 & areas\_verdes\_m2\_hab \\
top\_5\_menor\_areas\_verdes\_m2\_hab & 4.0 & fact\_desigualdad\_comunal & 13109 & LA CISTERNA & 0.904776 & areas\_verdes\_m2\_hab \\
top\_5\_menor\_areas\_verdes\_m2\_hab & 5.0 & fact\_desigualdad\_comunal & 13130 & SAN MIGUEL & 1.154347 & areas\_verdes\_m2\_hab \\
top\_5\_mayor\_pobreza\_ingresos\_pct & 1.0 & fact\_desigualdad\_comunal & 13112 & LA PINTANA & 9.2938 & pobreza\_ingresos\_pct \\
top\_5\_mayor\_pobreza\_ingresos\_pct & 2.0 & fact\_desigualdad\_comunal & 13131 & SAN RAMON & 6.8821 & pobreza\_ingresos\_pct \\
top\_5\_mayor\_pobreza\_ingresos\_pct & 3.0 & fact\_desigualdad\_comunal & 13116 & LO ESPEJO & 6.7817 & pobreza\_ingresos\_pct \\
top\_5\_mayor\_pobreza\_ingresos\_pct & 4.0 & fact\_desigualdad\_comunal & 13105 & EL BOSQUE & 6.1982 & pobreza\_ingresos\_pct \\
top\_5\_mayor\_pobreza\_ingresos\_pct & 5.0 & fact\_desigualdad\_comunal & 13104 & CONCHALI & 6.0799 & pobreza\_ingresos\_pct \\
top\_5\_menor\_ipp\_pesos\_hab & 1.0 & fact\_desigualdad\_comunal & 13103 & CERRO NAVIA & 27701.980354 & ipp\_pesos\_hab \\
top\_5\_menor\_ipp\_pesos\_hab & 2.0 & fact\_desigualdad\_comunal & 13112 & LA PINTANA & 38281.739358 & ipp\_pesos\_hab \\
top\_5\_menor\_ipp\_pesos\_hab & 3.0 & fact\_desigualdad\_comunal & 13117 & LO PRADO & 39312.805346 & ipp\_pesos\_hab \\
top\_5\_menor\_ipp\_pesos\_hab & 4.0 & fact\_desigualdad\_comunal & 13105 & EL BOSQUE & 45941.619379 & ipp\_pesos\_hab \\
top\_5\_menor\_ipp\_pesos\_hab & 5.0 & fact\_desigualdad\_comunal & 13111 & LA GRANJA & 47063.353627 & ipp\_pesos\_hab \\
\hline
\end{tabular}
\end{table}


En síntesis, la carga final reproduce el resultado del ETL sin alterar la centralidad de \texttt{codigo\_comuna}. \texttt{dim\_comuna} conserva los descriptores territoriales, \texttt{fact\_desigualdad\_comunal} concentra las variables finales cuantitativas y \texttt{metadata\_fuentes} mantiene la trazabilidad del proceso de lectura.
